\documentclass[11pt]{amsart}
\usepackage[margin=1in]{geometry}
\usepackage{amsmath,amssymb,amsthm,mathtools}
\usepackage{microtype}
\usepackage[hidelinks]{hyperref}
\usepackage{xcolor}
\usepackage{enumitem}
\usepackage{booktabs,array}

\newtheorem{theorem}{Theorem}[section]
\newtheorem{proposition}[theorem]{Proposition}
\newtheorem{lemma}[theorem]{Lemma}
\newtheorem{corollary}[theorem]{Corollary}
\theoremstyle{definition}
\newtheorem{remark}[theorem]{Remark}

\newcommand{\mot}{\operatorname{motion}}
\newcommand{\supp}{\operatorname{supp}}
\newcommand{\dist}{\operatorname{dist}}
\newcommand{\Aut}{\operatorname{Aut}}
\newcommand{\eps}{\varepsilon}
\newcommand{\gam}{\gamma}
\newcommand{\R}{\mathbb{R}}

\title[A cubic diameter bound for motion]{A cubic diameter bound for the motion\\of primitive distance-regular graphs}
\author{Machine-assisted proof draft}
\date{July 23, 2026}

\begin{document}

\begin{abstract}
Let $X$ be a primitive distance-regular graph on $n$ vertices, of diameter $d\ge3$.
We prove, modulo the published results quoted explicitly below, that either $X$ is a Johnson graph or a Hamming graph, or
\[
 \mot(X)\ge \frac{n}{14d^3}.
\]
The argument sharpens the diameter dependence in the primitive motion theorem of Pyber--Skresanov.  The principal new step is a cubic-scale Hamming stability lemma.  Its proof replaces the uniform additive loss in Kivva's standard-sequence argument by an exact relative-drop recurrence, and proves the required multiplicity surplus analytically; no finite parameter enumeration is used.  Additional ingredients are an exact formula for the vertices distinguishing an adjacent pair and a support-sensitive use of the geodesic boundary argument.
\end{abstract}
\maketitle

\begin{center}
\fbox{\parbox{0.93\textwidth}{\small
\textbf{Status.} This is a complete unrefereed proof draft, not a peer-reviewed result.  Every new implication used in the proof is written out; the dependency ledger in Section~\ref{sec:audit} lists the imported theorems and the hypotheses checked for each.  A targeted literature search through July 23, 2026 did not locate this explicit $d^{-3}$ bound, but no claim of priority is made without specialist review.}}
\end{center}

\section{Statement and notation}
Let $X$ be a primitive distance-regular graph on $n$ vertices, with valency $k$ and diameter $d\ge3$.  We use the standard intersection numbers
\[
 a_i=k-b_i-c_i,\qquad \lambda=a_1,\qquad \mu=c_2,
\]
and write $k_i$ for the size of a distance-$i$ sphere.  Let $\theta=\theta_1$ be the second largest adjacency eigenvalue and write the smallest eigenvalue as $-m$, where $m>0$.

\begin{theorem}\label{thm:main}
Either $X$ is a Johnson graph or a Hamming graph, or
\[
 \boxed{\mot(X)\ge \frac{n}{14d^3}.}
\]
\end{theorem}

Set
\begin{equation}\label{eq:parameters}
 \gam=\frac1{14d^3},\qquad
 \eps=\frac{2\gam}{1-\gam}=\frac2{14d^3-1},\qquad
 \alpha=\frac{1-\gam}{d}.
\end{equation}
We shall repeatedly use
\begin{equation}\label{eq:closure}
 \frac{\eps(1-\gam)}2=\gam,
 \qquad
 \eps>\gam,
 \qquad
 \eps<\frac17,
 \qquad
 \eps<0.0065.
\end{equation}

\section{Adjacent distinguishers and small supports}
For vertices $u,v$, define
\[
 D(u,v)=\{z:\dist(u,z)\ne\dist(v,z)\}.
\]
For adjacent $u,v$, the cardinality depends only on the distance relation; denote it by $D(1)$.

\begin{lemma}[Exact adjacent-pair identity]\label{lem:D1}
For every distance-regular graph of diameter at least three,
\begin{equation}\label{eq:D1}
 D(1)=2+\frac2k\sum_{i=2}^{d}k_i c_i.
\end{equation}
Consequently,
\begin{equation}\label{eq:D1mu}
 D(1)>\frac{\mu}{k}n.
\end{equation}
\end{lemma}

\begin{proof}
For adjacent $u,v$, the vertices not distinguishing them are those at a common distance $i$ from both.  Intersection-number balance gives
\[
 k p^1_{i,i}=k_i p^i_{1,i}=k_i a_i,
\]
so
\[
 D(1)=n-\frac1k\sum_{i=1}^{d}k_i a_i.
\]
Using $a_i=k-b_i-c_i$, $k_i b_i=k_{i+1}c_{i+1}$, $b_d=0$, and $k_1c_1=k$, we obtain
\[
 \sum_{i=1}^{d}k_i a_i=k(n-2)-2\sum_{i=2}^{d}k_i c_i,
\]
which proves \eqref{eq:D1}.

We also use the standard comparison $c_i\le b_j$ whenever $i+j\le d$.  Indeed, put $u,w,v$ in that order on a geodesic with $\dist(u,w)=i$ and $\dist(w,v)=j$.  Every neighbor $z$ of $w$ satisfying $\dist(u,z)=i-1$ has $\dist(v,z)=j+1$, so the $c_i$ such neighbors are among the $b_j$ such neighbors.  Taking $(i,j)=(2,1)$ gives $c_2\le b_1$, hence
\[
 k_2=\frac{kb_1}{c_2}\ge k.
\]
Thus $\sum_{i=2}^{d}k_i\ge(n-1)/2$.  Since $c_i\ge c_2=\mu$ for $i\ge2$ and $\mu\le k$, formula \eqref{eq:D1} yields
\[
 D(1)\ge2+\frac{\mu}{k}(n-1)>\frac{\mu}{k}n.
\]
\end{proof}

\begin{lemma}[Support-sensitive boundary]\label{lem:boundary}
Let $g\in\Aut(X)$ be nonidentity, let $S=\supp(g)$, and put $\rho=|S|/n\le1/2$.  Some $x\in S$ has at least
\begin{equation}\label{eq:fixedneighbors}
 \frac{k}{d}(1-\rho)
\end{equation}
fixed neighbors.
\end{lemma}

\begin{proof}
We spell out the final, support-sensitive form of the geodesic-load argument in \cite[Proposition~2.8]{PS}.  Let
\[
 \partial^+(S)=\{(u,v):u\in S,\ v\notin S,\ u\sim v\}
\]
be the outgoing boundary in the set of oriented edges.  For ordered vertices $a,b$, let $p(a,b)$ be the number of geodesics from $a$ to $b$.  For an oriented edge $e$, let $N_e(a,b)$ count those geodesics traversing $e$ in its given orientation.  The coherent-configuration intersection-number identities show that
\[
 P_e=\sum_{a,b}\frac{N_e(a,b)}{p(a,b)}
\]
is independent of $e$; write the common value as $P$.  There are $nk$ oriented edges, hence
\[
 nkP=\sum_{a,b}\dist(a,b)\le dn^2,
 \qquad P\le\frac{dn}{k}.
\]
Every geodesic from $a\in S$ to $b\notin S$ crosses an edge of $\partial^+(S)$, so
\[
 |S|(n-|S|)\le |\partial^+(S)|P.
\]
Consequently,
\[
 |\partial^+(S)|\ge |S|\frac{k}{d}\frac{n-|S|}{n}
 =|S|\frac{k}{d}(1-\rho).
\]
Averaging over the initial vertices in $S$ proves the claim.  Notice that all edges here are oriented, so no factor-of-two convention is hidden.
\end{proof}

\begin{lemma}[A small support moves an adjacent pair]\label{lem:adjacent}
Assume the notation of Lemma~\ref{lem:boundary}.  If
\[
 \mu<\frac{k}{d}(1-\rho),
\]
then some $x\in S$ satisfies $x\sim x^g$, and
\begin{equation}\label{eq:supportdata}
 \lambda\ge\frac{k}{d}(1-\rho),
 \qquad
 D(1)\le |S|.
\end{equation}
\end{lemma}

\begin{proof}
Choose $x$ as in Lemma~\ref{lem:boundary}.  Every fixed neighbor of $x$ is also adjacent to $x^g$.  If $x$ and $x^g$ were nonadjacent, then they would be at distance two and would have exactly $\mu$ common neighbors, contradicting the displayed hypothesis.  Hence $x\sim x^g$, and the fixed common neighbors give the bound on $\lambda$.

For every $z\notin S$,
\[
 \dist(x,z)=\dist(x^g,z^g)=\dist(x^g,z),
\]
so no fixed vertex distinguishes $x$ from $x^g$.  Thus $D(1)\le|S|$.
\end{proof}

\section{Small support forces Delsarte geometry}
The proof of \cite[Proposition~2.6]{PS}, retaining the full expression furnished by Metsch's theorem, shows that a sub-amply regular graph with $\lambda^2\ge4k\mu$ contains a clique of size at least
\begin{equation}\label{eq:metsch}
 \lambda+2-
 \left(\left\lceil\frac{3k}{2(\lambda+1)}\right\rceil-1\right)(\mu-1).
\end{equation}

\begin{proposition}\label{prop:structure}
Suppose a nonidentity automorphism of $X$ has support density $\rho<\gam$.  Then
\begin{equation}\label{eq:structdata}
 \mu<\gam k,
 \qquad
 \lambda>\alpha k,
 \qquad
 k>14d^3,
\end{equation}
and $X$ is Delsarte-geometric with smallest eigenvalue $-m$ satisfying
\begin{equation}\label{eq:mleqd}
 m\le d.
\end{equation}
\end{proposition}

\begin{proof}
We first show that $\mu<k/(2d)$.  Suppose instead that $\mu\ge k/(2d)$.  Let $k_i=k_{\max}$ be the largest nontrivial relation valency.  Since $k_2\ge k_1$, we may choose $i\ge2$.  Monotonicity of the $c_i$ gives
\[
 \frac{k_i}{k_{i-1}}=\frac{b_{i-1}}{c_i}\le\frac{k}{\mu}\le2d,
\]
so $n-k_{\max}\ge k_{i-1}\ge k_{\max}/(2d)$.  Propositions~2.10 and~2.12 of \cite{PS} give
\[
 \mot(X)\ge\frac{n-k_{\max}}d.
\]
If $k_{\max}\ge n/2$, this is at least $n/(4d^2)$; otherwise it is greater than $n/(2d)$.  Either bound exceeds $\gam n$, contradicting $\rho<\gam$.

Since $\rho<\gam<1/2$, Lemma~\ref{lem:adjacent} now applies and gives
\[
 \lambda>\frac{1-\gam}{d}k=\alpha k.
\]
Together with Lemma~\ref{lem:D1} and $D(1)\le|S|$, it gives
\[
 \mu<\rho k<\gam k.
\]
As $\mu\ge1$, we have $k>1/\gam=14d^3$.

The following three scalar inequalities hold for $d\ge3$:
\begin{align}
 \alpha^2&>4\gam,\label{eq:scalar1}\\
 \alpha-\frac{3\gam}{2\alpha}&>\frac1{d+1},\label{eq:scalar2}\\
 \alpha&>(d+1)^2\gam.\label{eq:scalar3}
\end{align}
Their exact positive numerators are recorded in Appendix~\ref{sec:scalar}.  Inequality \eqref{eq:scalar1} gives $\lambda^2>4k\mu$, so \eqref{eq:metsch} applies.  If $L$ is the resulting clique size, then $\lceil x\rceil-1<x$ and \eqref{eq:scalar2} give
\begin{align*}
 L-1
 &\ge \lambda+1-
 \left(\left\lceil\frac{3k}{2(\lambda+1)}\right\rceil-1\right)(\mu-1)\\
 &>k\left(\alpha-\frac{3\gam}{2\alpha}\right)
 >\frac{k}{d+1}.
\end{align*}
The Delsarte clique bound $L\le1+k/m$ implies $m<d+1$.  Finally, \eqref{eq:scalar3} gives
\[
 m^2\mu<(d+1)^2\gam k<\alpha k<\lambda.
\]
The Bang--Koolen criterion \cite[Proposition~2.5]{PS} makes $X$ Delsarte-geometric.  In a Delsarte geometry $m$ is an integer by \cite[Lemma~2.3]{PS}; hence $m\le d$.
\end{proof}

\section{A dominant distance and a high second eigenvalue}
Assume for the remainder of the proof that a nonidentity automorphism $g$ has support density $\rho<\gam$.  Proposition~\ref{prop:structure} is therefore available.

\begin{lemma}[Transition without a diameter loss]\label{lem:transition}
If
\[
 b_j\ge\eps k,
 \qquad
 c_{j+1}\ge\eps k
\]
for some $1\le j\le d-1$, then $|\supp(g)|>\eps n>\gam n$, a contradiction.
\end{lemma}

\begin{proof}
Monotonicity gives $b_i\ge\eps k$ for $i\le j$ and $c_i\ge\eps k$ for $i\ge j+1$.  Thus $a_i\le(1-\eps)k$ for every $i\ge1$, and
\[
 D(1)=n-\frac1k\sum_{i=1}^{d}k_i a_i
 \ge n-(1-\eps)(n-1)>\eps n.
\]
Use $D(1)\le|\supp(g)|$ from Lemma~\ref{lem:adjacent} and $\eps>\gam$.
\end{proof}

Let $t$ be the least index for which $b_t\le\eps k$; it exists since $b_d=0$.  The standard inequality $2\lambda\le k+\mu$ gives
\[
 b_1=k-\lambda-1\ge\frac{k-\mu-2}{2}
 >\frac{1-3\gam}{2}k>\frac{k}{3}>\eps k.
\]
Here we used $\mu<\gam k$ and $1/k<\gam$.  Thus $t\ge2$.  Lemma~\ref{lem:transition}, applied at $j=t-1$, gives
\begin{equation}\label{eq:dominant}
 b_t\le\eps k,
 \qquad
 c_t<\eps k.
\end{equation}

\begin{lemma}[High second eigenvalue]\label{lem:hightheta}
Under these assumptions,
\[
 \theta\ge(1-\eps)b_1.
\]
\end{lemma}

\begin{proof}
Suppose $\theta<(1-\eps)b_1$.  Since $m\le d$, $k>14d^3$, and $b_1>k/3$, we also have $m<(1-\eps)b_1$.  Hence the zero-weight spectral radius $\xi=\max\{\theta,m\}$ is less than $(1-\eps)b_1$.

Moreover, $\lambda>\mu$, since $\alpha>\gam$.  Every pair of vertices has at most $\lambda$ common neighbors.  Babai's spectral motion bound \cite[Proposition~2.13]{PS} therefore gives
\begin{align*}
 \rho
 &\ge\frac{k-\xi-\lambda}{k}
 >\frac{1+\eps b_1}{k}\\
 &=\frac{1-\eps}{k}+\eps\frac{k-\lambda}{k}
 \ge\frac{1-\eps}{k}+\frac\eps2\left(1-\frac\mu k\right)\\
 &>\frac\eps2(1-\gam)=\gam,
\end{align*}
contradicting the choice of $g$.  The penultimate inequality uses $2\lambda\le k+\mu$, and the final identity is \eqref{eq:closure}.
\end{proof}

\section{The case $\mu\ge3$}
\begin{proposition}\label{prop:johnson}
If $\mu\ge3$, then $X$ is a Johnson graph.
\end{proposition}

\begin{proof}
If the neighborhood graphs are disconnected, Proposition~2.20 of \cite{PS} gives
\[
 \theta+1\le\frac57b_1,
\]
contradicting Lemma~\ref{lem:hightheta} because $\eps<2/7$.  Hence a neighborhood graph is connected.

Proposition~2.19 of \cite{PS}, which quotes Kivva's Johnson characterization, supplies an absolute constant $\eps_*>0.0065$.  Its hypotheses hold here: $\mu\ge2$; Lemma~\ref{lem:hightheta} and $\eps<0.0065<\eps_*$ give
\[
 \theta+1>(1-\eps_*)b_1;
\]
and $k>14d^3\ge14m^3>\max\{m^3,29\}$.  Therefore $X$ is Johnson.
\end{proof}

\section{A cubic-scale Hamming stability lemma}\label{sec:hamming}
This section supplies the step absent from the earlier drafts.

\begin{proposition}[Cubic-scale Hamming stability]\label{prop:hamming}
Let $X$ be a Delsarte-geometric distance-regular graph of diameter $d\ge2$, with smallest eigenvalue $-m$.  Suppose
\[
 \mu=2,
 \qquad
 m\le d,
 \qquad
 0<\eps=\frac2{14d^3-1},
\]
and for some $2\le t\le d$,
\[
 b_t\le\eps k,
 \qquad
 c_t<\eps k,
 \qquad
 \theta\ge(1-\eps)b_1.
\]
Then $X$ is a Hamming graph.
\end{proposition}

\subsection{Geometric parameters and the standard sequence}
Let $\tau_i,\psi_i$ be the geometric parameters of Kivva \cite[Lemma~2.16]{Kivva}.  Thus
\begin{equation}\label{eq:geometricparams}
 c_i=\tau_i\psi_{i-1},
 \qquad
 b_i=(m-\tau_i)\left(\frac{k}{m}+1-\psi_i\right).
\end{equation}
Since $\mu=\tau_2\psi_1=2$ and $\tau_2\ge\psi_1$, we have
\begin{equation}\label{eq:mu2params}
 \tau_2=2,
 \qquad
 \psi_1=1,
 \qquad
 b_1=\frac{m-1}{m}k.
\end{equation}
By \cite[Lemma~2.19]{Kivva}, every neighborhood graph is a disjoint union of $m$ cliques.

The bound $m\eps<1/60$ follows from $m\le d$ and
\[
 \frac{2d}{14d^3-1}<\frac1{60}
 \qquad(d\ge3).
\]
In particular, $\eps<1/m^2$.  Kivva's strict-growth lemma \cite[Lemma~4.2]{Kivva} applies and gives
\begin{equation}\label{eq:taugrowth}
 \tau_i<\tau_{i+1}\qquad(1\le i\le t-2).
\end{equation}
Since $b_{t-1}>0$, one has $\tau_{t-1}\le m-1$.  Consequently
\begin{equation}\label{eq:trange}
 2\le t\le m,
 \qquad
 r:=m-\tau_{t-1}\in\{1,\ldots,m-t+1\}.
\end{equation}

Let $(u_i)_{i=0}^{d}$ be the standard sequence corresponding to $\theta$.  Put
\[
 y_i=\frac{u_{i-1}-u_i}{u_{i-1}}
 \qquad(i\ge1)
\]
whenever the preceding terms are positive.  The standard-sequence recurrence is exactly equivalent to
\begin{equation}\label{eq:riccati}
 y_{i+1}=
 \frac{k-\theta+c_i\,y_i/(1-y_i)}{b_i}.
\end{equation}
Define
\begin{equation}\label{eq:ABC}
 B=1+(m-1)\eps,
 \qquad
 C=1+\left(\frac53m-1\right)\eps,
 \qquad
 A=\frac{C}{1-m\eps}=1+\delta.
\end{equation}
The elementary bounds $m\eps<1/60$ imply
\begin{equation}\label{eq:Adelta}
 A<\frac65,
 \qquad
 \delta=\frac{(8m-3)\eps}{3(1-m\eps)}<3m\eps.
\end{equation}

\begin{lemma}[Relative-drop estimate]\label{lem:drops}
For $2\le i\le t-1$,
\begin{equation}\label{eq:ybound}
 0<y_i\le\frac{A}{m-\tau_{i-1}}.
\end{equation}
Consequently, with $H_0=0$ and $H_j=\sum_{h=1}^{j}1/h$,
\begin{equation}\label{eq:ulower}
 u_{t-1}\ge
 u_1\frac{r}{r+t-2}\left(1-\delta H_{t-2}\right).
\end{equation}
\end{lemma}

\begin{proof}
From \eqref{eq:mu2params} and $\theta\ge(1-\eps)b_1$,
\begin{equation}\label{eq:ktheta}
 k-\theta\le\frac{B}{m}k,
 \qquad
 u_1=\frac\theta k\ge(1-\eps)\frac{m-1}{m}>\frac37.
\end{equation}
Thus $0<y_1/(1-y_1)<4/3$.  Since $c_t\ge c_2=2$ and $c_t<\eps k$, we also have $1/k<\eps/2$.  Applying \eqref{eq:riccati} at $i=1$, using $c_1=1$ and $b_1=(m-1)k/m$, gives
\[
 y_2<\frac{B+(2m/3)\eps}{m-1}
 =\frac{C}{m-1}<\frac{A}{m-1}.
\]
This proves the base case whenever $t\ge3$; for $t=2$ there is nothing to prove.

For $i\le t-1$, monotonicity and \eqref{eq:geometricparams} give
\[
 c_i\le c_t<\eps k,
 \qquad
 b_i\ge(m-\tau_i)\left(\frac1m-\eps\right)k.
\]
Suppose $2\le i\le t-2$ and \eqref{eq:ybound} holds.  By strict growth,
$\tau_{i-1}\le m-3$, so $y_i\le A/3<2/5$ and hence
$y_i/(1-y_i)<2/3$.  Formula \eqref{eq:riccati} then yields
\[
 y_{i+1}
 \le\frac{B+(2m/3)\eps}{(m-\tau_i)(1-m\eps)}
 =\frac{A}{m-\tau_i}.
\]
The induction also proves positivity of all terms involved.

By \eqref{eq:taugrowth}, for $2\le i\le t-1$,
\[
 m-\tau_{i-1}\ge r+t-i.
\]
Therefore
\begin{align*}
 \frac{u_{t-1}}{u_1}
 &=\prod_{i=2}^{t-1}(1-y_i)\\
 &\ge\prod_{s=r+1}^{r+t-2}\left(1-\frac{1+\delta}{s}\right)\\
 &=\frac{r}{r+t-2}
   \prod_{j=r}^{r+t-3}\left(1-\frac{\delta}{j}\right)\\
 &\ge\frac{r}{r+t-2}\left(1-\delta H_{t-2}\right).
\end{align*}
The last inequality is the elementary product bound
$\prod(1-a_j)\ge1-\sum a_j$ for $0\le a_j\le1$.
\end{proof}

\subsection{Concentration at the dominant sphere}
Set
\begin{equation}\label{eq:q}
 q=\frac{m\eps}{1-m\eps}<1.
\end{equation}
For $2\le i\le t$,
\[
 \frac{k_{i-1}}{k_i}=\frac{c_i}{b_{i-1}}\le q,
\]
and the same inequality holds for $i=1$ because $k_0/k_1=1/k<\eps/2<q$.  For $t\le i\le d-1$, the inequality $b_i\le b_t\le\eps k$ and formula \eqref{eq:geometricparams} give
\[
 \psi_i\ge\left(\frac1m-\eps\right)k,
 \qquad
 c_{i+1}\ge\psi_i,
 \qquad
 \frac{k_{i+1}}{k_i}=\frac{b_i}{c_{i+1}}\le q.
\]
Summing both geometric tails around $k_t$ yields
\begin{equation}\label{eq:ktmass}
 n\le k_t\left(1+2\frac{q}{1-q}\right),
 \qquad
 k_t\ge\frac{1-q}{1+q}n=(1-2m\eps)n.
\end{equation}

\subsection{The multiplicity surplus}
Since $\psi_{t-1}\ge1$, formula \eqref{eq:geometricparams} gives
\[
 b_{t-1}\le\frac{r}{m}k.
\]
Using $k_{t-1}b_{t-1}=k_t c_t$, \eqref{eq:ulower}, \eqref{eq:ktmass}, and
$u_1\ge(1-\eps)(m-1)/m$, we obtain
\begin{equation}\label{eq:FR}
 k_{t-1}u_{t-1}^2\ge\frac nk\,F R,
\end{equation}
where
\begin{align}
 F&=(1-2m\eps)(1-\eps)^2
      \left(1-\delta H_{t-2}\right)^2,
      \label{eq:F}\\
 R&=\frac{c_t r(m-1)^2}{m(r+t-2)^2}.
      \label{eq:R}
\end{align}

\begin{lemma}[Exact surplus factor]\label{lem:R}
Unless
\begin{equation}\label{eq:endpoint}
 c_t=t=m=d,
\end{equation}
one has
\begin{equation}\label{eq:Rbound}
 R\ge1+\frac1m.
\end{equation}
\end{lemma}

\begin{proof}
By \eqref{eq:taugrowth}, $\tau_{t-1}\ge t-1$; since $\psi_{t-2}\ge1$ and the $c_i$ are nondecreasing,
\begin{equation}\label{eq:ctlower}
 c_t\ge c_{t-1}=\tau_{t-1}\psi_{t-2}\ge t-1.
\end{equation}

Suppose first that $c_t=t-1$.  Equality holds throughout \eqref{eq:ctlower}, so
$\tau_{t-1}=t-1$ and $r=m-t+1$.  One has $t\ge4$: the cases $t=2$ and $t=3$ contradict respectively $c_2=2$ and $c_3>c_2$ \cite[Corollary~2.8]{Kivva}.  Moreover, $c_t=c_{t-1}$.  The induced-quadrangle form of Terwilliger's inequality, used in \cite[proof of Proposition~4.6]{Kivva}, gives
\[
 b_{t-1}\ge b_t+\lambda+2\ge\frac{k}{m}+1.
\]
But $b_{t-1}\le r k/m$, so $r\ge2$, equivalently $t\le m-1$.  Hence $4\le t\le m-1$ and
\[
 R=\frac{(t-1)(m-t+1)}m\ge\frac{m+1}{m}.
\]
The final inequality follows because the concave product $(t-1)(m-t+1)$ has its minimum on this interval at an endpoint, where it is at least $m+1$.

Now suppose $c_t\ge t$.  If $t\le m-1$, then it suffices to prove
\begin{equation}\label{eq:rineq}
 \frac{t r(m-1)^2}{(r+t-2)^2}\ge m+1
 \qquad(1\le r\le m-t+1).
\end{equation}
The function $r/(r+t-2)^2$ has no interior minimum, so it is enough to check the two endpoints.  At $r=1$, \eqref{eq:rineq} is
\[
 t(m-1)^2\ge(m+1)(t-1)^2;
\]
for fixed $t$ the difference is increasing in $m\ge t+1$, and at $m=t+1$ it equals $3t-2>0$.  At $r=m-t+1$, it becomes
\[
 t(m-t+1)\ge m+1,
\]
which is equivalent to $(t-1)(m-t)\ge1$.

It remains to consider $t=m$.  Then $r=1$ and $R=c_t/m$.  Thus \eqref{eq:Rbound} holds if $c_t\ge m+1$.  If $c_t=m=t$, we claim $t=d$.  If $t<d$, then $b_t\ge1$, while $\tau_{t-1}=m-1$ and the induced-quadrangle Terwilliger inequality give
\[
 b_{t-1}\ge c_{t-1}-c_t+b_t+\lambda+2\ge\lambda+2\ge\frac{k}{m}+1.
\]
On the other hand, \eqref{eq:geometricparams} gives $b_{t-1}\le k/m$, a contradiction.  Hence $c_t=t=m=d$, which is precisely \eqref{eq:endpoint}.
\end{proof}

\begin{lemma}[The analytic loss is smaller than the surplus]\label{lem:F}
One has
\begin{equation}\label{eq:Fbound}
 F>\frac{m}{m+1}.
\end{equation}
\end{lemma}

\begin{proof}
From \eqref{eq:Adelta}, $H_{t-2}\le d-2$, and $m\le d$,
\begin{align*}
 F
 &\ge1-2(m+1)\eps-2\delta H_{t-2}\\
 &>1-\eps(6d^2-10d+2).
\end{align*}
The scalar inequality
\begin{equation}\label{eq:lossineq}
 \eps(6d^2-10d+2)<\frac1{d+1}
\end{equation}
has positive cross-multiplied difference
\[
 2d^3+8d^2+16d-5>0.
\]
Therefore
\[
 F>\frac{d}{d+1}\ge\frac{m}{m+1}.
\]
\end{proof}

\begin{proof}[Proof of Proposition~\ref{prop:hamming}]
If the endpoint \eqref{eq:endpoint} does not hold, Lemmas~\ref{lem:R} and~\ref{lem:F} give $FR>1$.  By \eqref{eq:FR},
\[
 k_{t-1}u_{t-1}^2>\frac nk.
\]
Biggs' multiplicity formula \cite[Theorem~2.10]{Kivva} then yields
\[
 f_1=\frac{n}{\sum_{i=0}^{d}k_i u_i^2}<k.
\]
Terwilliger's local-eigenvalue theorem \cite[Theorem~4.1]{Kivva} says that every neighborhood graph has the eigenvalue
\[
 -1-\frac{b_1}{\theta+1}<-1.
\]
This is impossible because every neighborhood graph is a disjoint union of cliques, whose least eigenvalue is $-1$.

Hence $c_d=m=d$.  Corollary~4.3 of \cite{Kivva} and the integrality of the $\tau_i$ give $\tau_i=i$ for all $i$.  Since
$c_d=\tau_d\psi_{d-1}=d$, one has $\psi_{d-1}=1$.  If some $\psi_{i-1}\ge2$ were followed by $\psi_i=1$, then
\[
 i+1=c_{i+1}=\tau_{i+1}\psi_i
 \ge c_i=\tau_i\psi_{i-1}\ge2i,
\]
a contradiction.  Thus $\psi_i=1$ for every $i$, and \eqref{eq:geometricparams} gives
\[
 c_i=i,
 \qquad
 b_i=(d-i)\frac{k}{d}.
\]
The intersection array is that of $H(d,1+k/d)$.  Egawa's characterization \cite[Theorem~2.24]{Kivva} leaves only a Hamming graph or a Doob graph.  A Doob graph would require $1+k/d=4$, i.e. $k=3d$, contrary to $k>14d^3$.  Hence $X$ is Hamming.
\end{proof}

\section{The case $\mu=1$}
\begin{proposition}\label{prop:mu1}
If $\mu=1$, then $\mot(X)>\gam n$.
\end{proposition}

\begin{proof}
First suppose $m\ge3$.  Let $\widetilde X$ be the dual graph of the Delsarte geometry.  Kivva's Lemma~2.26 gives
\begin{equation}\label{eq:dualdegree}
 \widetilde k=(m-1)\left(1+\frac{k}{m}\right),
\end{equation}
and the discussion following \cite[Lemma~5.1]{Kivva} shows that every pair of vertices of $\widetilde X$ has at most
\[
 q=m-2
\]
common neighbors.  Since $k>14d^3\ge m^2$, Lemma~2.27 of \cite{Kivva} applies to the spectrum of $\widetilde X$.

Put $\eta=1/(8d^2)$.  Proposition~2.9 of \cite{PS} gives
\[
 \theta\le(1-\eta)k.
\]
For every nontrivial eigenvalue $\widetilde\theta$ of $\widetilde X$, the spectral inclusion in Kivva's lemma gives
\[
 \widetilde\theta\le\widetilde k-\eta k\le(1-\eta)\widetilde k,
\]
because $k\ge\widetilde k$.  On the negative side,
\[
 \widetilde\theta\ge-\frac{k}{m}-1=-\frac{\widetilde k}{m-1}
 \ge-(1-\eta)\widetilde k,
\]
since $m\ge3$ and $\eta<1/2$.  Thus the zero-weight spectral radius of $\widetilde X$ is at most $(1-\eta)\widetilde k$.

Moreover,
\[
 \frac q{\widetilde k}<\frac{m}{k}\le\frac d{k}<\frac1{14d^2}.
\]
Babai's spectral motion bound therefore gives
\[
 \frac{\mot(\widetilde X)}{|V(\widetilde X)|}
 >\frac1{8d^2}-\frac1{14d^2}
 =\frac3{56d^2}.
\]
Kivva's motion-transfer Corollary~5.6 now yields
\[
 \mot(X)>\frac{3n}{112d^2}>\frac{n}{14d^3}.
\]

If $m=2$, Proposition~5.13 of \cite{Kivva} gives $\mot(X)\ge n/16>\gam n$.  The case $m=1$ cannot occur: every vertex would lie in exactly one clique of the Delsarte geometry, so connectedness would force $X$ itself to be complete, contrary to $d\ge3$.
\end{proof}

\section{Proof of the main theorem}
\begin{proof}[Proof of Theorem~\ref{thm:main}]
Suppose, toward a contradiction, that $X$ is neither Johnson nor Hamming and that a nonidentity automorphism $g$ has support density $\rho<\gam$.  Proposition~\ref{prop:structure} gives Delsarte geometry and $m\le d$.  Lemma~\ref{lem:transition} produces an index $t$ satisfying \eqref{eq:dominant}, and Lemma~\ref{lem:hightheta} gives $\theta\ge(1-\eps)b_1$.

If $\mu\ge3$, Proposition~\ref{prop:johnson} makes $X$ Johnson.  If $\mu=2$, Proposition~\ref{prop:hamming} makes $X$ Hamming.  If $\mu=1$, Proposition~\ref{prop:mu1} gives $\mot(X)>\gam n$.  Every case contradicts the assumptions.  Hence every nonexceptional graph has motion at least $\gam n=n/(14d^3)$.
\end{proof}

\section{Dependency and hypothesis audit}\label{sec:audit}
The proof imports the following published statements.

\begin{center}
\small
\begin{tabular}{@{}>{\raggedright\arraybackslash}p{0.25\textwidth}p{0.65\textwidth}@{}}
\toprule
Imported result & Hypotheses checked in this draft \\
\midrule
PS Proposition 2.8 & Used only through its coherent-configuration geodesic-load uniformity; the stronger support-sensitive final inequality is rederived in Lemma~\ref{lem:boundary} with oriented edges. \\
PS Propositions 2.10, 2.12 & $X$ is primitive, and every nontrivial relation in its distance scheme has diameter at most $d$. \\
Metsch expression in PS Proposition 2.6 & $\lambda^2>4k\mu$ follows from \eqref{eq:scalar1}; the ceiling term is retained rather than replaced by $\lambda/2$. \\
Delsarte and Bang--Koolen & The clique first gives $m<d+1$; then \eqref{eq:scalar3} gives $m^2\mu<\lambda$. \\
PS Proposition 2.13 & The maximum common-neighbor count is $\lambda$ because $\lambda>\mu$; the zero-weight spectral radius is controlled explicitly. \\
PS Propositions 2.19--2.20 & $d>2$, $\mu\ge3$, Delsarte geometry, connected/disconnected neighborhoods, $\eps<0.0065$, and $k\ge\max\{m^3,29\}$. \\
Kivva Lemmas 2.16--2.19, 4.2 & Delsarte geometry, $\mu=2$, $c_t<\eps k$, and $\eps<1/m^2$. \\
Biggs and Terwilliger & The new estimates prove $f_1<k$ before the local-eigenvalue theorem is invoked; local graphs are disjoint unions of cliques. \\
Egawa/Kivva Theorem 2.24 & The endpoint intersection array is proved explicitly; $k>14d^3$ excludes the Doob parameter $k=3d$. \\
Kivva Lemmas 2.26--2.27 and Corollary 5.6 & In the $\mu=1$ branch, $k>14d^3\ge m^2$, $m\ge3$, the dual spectral radius and common-neighbor count are bounded explicitly, and motion is transferred with the published factor $1/2$. \\
Kivva Proposition 5.13 & Used only for $\mu=1$, $m=2$, and $k>4$. \\
\bottomrule
\end{tabular}
\end{center}

\begin{remark}
The logical gap in the withdrawn $d^{-3}$ drafts was the unproved multiplicity surplus in the $\mu=2$ branch.  Lemma~\ref{lem:R} closes that gap: Kivva's published argument records a coarser lower bound in one subcase, but the complete admissible constraints imply the sharper uniform estimate $R\ge1+1/m$.  The recurrence and loss estimates then prove $FR>1$ symbolically, with no tuple enumeration.
\end{remark}

\appendix
\section{Scalar certificates}\label{sec:scalar}
For $\gam,\alpha$ from \eqref{eq:parameters}, the three inequalities used in Proposition~\ref{prop:structure} reduce to positivity of
\begin{align*}
 196d^6-56d^5-28d^3+1,
 \qquad
 175d^6-21d^5-14d^4-28d^3+d+1,
 \qquad
 13d^3-2d^2-d-1.
\end{align*}
After writing $d=x+3$, these become respectively
\begin{align*}
 &196x^6+3472x^5+25620x^4+100772x^3
   +222768x^2+262332x+128521,\\
 &175x^6+3129x^5+23296x^4+92414x^3
   +205947x^2+244378x+120586,\\
 &13x^3+115x^2+338x+329,
\end{align*}
which are positive for $x\ge0$.

The loss inequality \eqref{eq:lossineq} is equivalent to
\[
 2d^3+8d^2+16d-5>0;
\]
after $d=x+3$, the left side is
\[
 2x^3+26x^2+118x+169.
\]
Finally,
\[
 m\eps\le d\eps=\frac{2d}{14d^3-1}<\frac1{60}
\]
is equivalent to $14d^3-120d-1>0$, which holds at $d=3$ and is increasing thereafter.

\begin{thebibliography}{9}
\bibitem{PS}
L. Pyber and S. V. Skresanov,
\emph{On the automorphism group of a distance-regular graph},
J. Combin. Theory Ser. B \textbf{172} (2025), 94--114;
\href{https://arxiv.org/abs/2312.00383}{arXiv:2312.00383}.

\bibitem{Kivva}
B. Kivva,
\emph{A characterization of Johnson and Hamming graphs and proof of Babai's conjecture},
J. Combin. Theory Ser. B \textbf{151} (2021), 339--374;
\href{https://arxiv.org/abs/1912.11427}{arXiv:1912.11427}.
\end{thebibliography}

\end{document}
